\section*{Overview}

Topological sort turns a DAG into a linear order that respects every directed edge. Once that order exists, many
dependency problems become ordinary left-to-right DP or greedy scans.

\section*{When to Use It}

Use topological sort when:

\begin{itemize}
  \item the graph is directed and acyclic, or the task asks you to detect whether it is,
  \item edges represent prerequisites or dependency constraints,
  \item a DP transition only depends on earlier states.
\end{itemize}

\section*{Core Idea}

In a DAG, at least one vertex has indegree zero. Remove such vertices one by one, appending them to the order and
updating the indegrees of their outgoing neighbors. That is Kahn's algorithm.

\section*{Key Insight}

The order is not usually unique, and that is fine. What matters is that every edge goes from earlier to later. That is
exactly the condition needed to process a node after all of its prerequisites are done.

\section*{Operations / Main Technique}

\begin{itemize}
  \item compute indegrees,
  \item push all indegree-zero vertices into a queue,
  \item repeatedly pop, output, and relax outgoing edges,
  \item if not all vertices are output, a cycle exists.
\end{itemize}

\paragraph{Worked example.}
If course \(A\) and \(B\) must be taken before \(C\), then edges \(A \to C\) and \(B \to C\) guarantee \(C\) appears
after both of them in every topological order.

\section*{Correctness Intuition}

A vertex with indegree zero has no remaining prerequisite, so placing it next cannot violate any edge constraint. When
it is removed, its outgoing edges are no longer relevant to the remaining graph. Repeating this preserves correctness.

\section*{Complexity Analysis}

Kahn's algorithm runs in \(O(n + m)\) time and \(O(n)\) extra memory.

\section*{Implementation}

The sample returns an empty vector if the graph contains a cycle. That makes it easy to use the same routine both for
ordering and for cycle detection.

\section*{Common Pitfalls}

\begin{itemize}
  \item Forgetting that topological order exists only for DAGs.
  \item Using an order-based DP without checking for cycles first.
  \item Mixing edge direction with the dependency statement from the problem.
\end{itemize}

\section*{Variants / Extensions}

\begin{itemize}
  \item DFS finishing-order topological sort.
  \item Lexicographically smallest topological order with a priority queue.
  \item DAG DP once the order is available.
\end{itemize}

\section*{Practice Problems}

\begin{itemize}
  \item Order courses by prerequisites.
  \item Longest path in a DAG.
  \item Count paths in a DAG modulo a prime.
\end{itemize}

\section*{References}

\begin{itemize}
  \item \href{https://cp-algorithms.com/graph/topological-sort.html}{cp-algorithms: Topological Sorting}.
  \item \href{https://usaco.guide/gold/toposort?lang=cpp}{USACO Guide: Topological Sort}.
  \item \href{https://codeforces.com/catalog}{Codeforces Catalog}.
\end{itemize}
